% 本文件由 LaTeX 排版 Agent 根据有效 translation.json 自动生成。
% author=Origen; work_id=1014; title=致国瑞里书

\chapter{致\Person{国瑞里}书}

% structure: p0001 source=h1 target=omitted reason=page-title-already-rendered-by-parent

论从哲学所得之学何以有助于解释圣经；附热切的个人恳请。\par

此函致\Person{国瑞里}——后为\Place{凯撒勒雅}主教，号显灵者——原载\WorkTitle{斐罗卡利亚}（即\Person{尼撒的国瑞里}与\Person{凯撒勒雅的巴西略}所辑\Person{奥利振}著作选集）。\Person{德拉吕}与\Person{洛马茨}将其置于著作集卷首。此函为注释书之佳序，显示\Person{奥利振}视圣经研究为诸学之冠，且视科学知识——彼本人即为其中大师——仅为这至高学问之预备。\Person{德雷泽克}证明此函写于约235年，时\Person{奥利振}在\Place{凯撒勒雅}教授\Person{国瑞里}数年后，因\Person{马克西米努斯·特拉克斯}迫害而逃往\Place{卡帕多细亚}；而\Person{国瑞里}，据函中语气，已赴\Place{埃及}。\Person{国瑞里}于约239年在\Place{凯撒勒雅}——彼时学校于迫害后重聚——所发表的\WorkTitle{颂奥利振辞}，显示导师对其弟子真正进步的关切并未落空。\par

% structure: p0004 source=h2 target=section reason=relative-heading-rank; base=h2

\section{一、\Person{国瑞里}受劝将外教学问用于研读圣经}

愿您在天主内万安，最卓越可敬的先生，\Person{国瑞里}子，自\Person{奥利振}。天生的敏悟，如您深知，若勤加操练，适于产生一作品，使其拥有者——容我如此说——尽可能达至其所欲实践之艺的圆满；而您的天赋已足以使您成为一位卓越的罗马法学家，亦为最著名学派的希腊哲学家。但我对您的期望是：您应将自己的全部智力导向以基督信仰为终向，并以生产的方式。我且愿您一方面取希腊哲学中那些适宜作为基督信仰之普通或预备学科的部分，另一方面取几何学与天文学中足以助益解释圣经的部分。哲学家们谈论几何、音乐、文法、修辞与天文学为哲学的辅助；同样，我们也可说哲学本身为基督信仰的辅助。\par

% structure: p0006 source=h2 target=section reason=relative-heading-rank; base=h2

\section{二、此程序由夺取埃及人财物之故事预表}

或许这正是\BibleBook{出谷}{纪}所记天主指示以色列子民时隐晦表明的某种事。他们受命向邻舍及住在他们帐幕中的人求取银器、金器及衣裳；如此他们夺取埃及人，获得材料以制造奉命为敬拜天主所备之物。因为以色列子民从埃及人夺来的东西中，制造了至圣所的器物：约柜及其盖、革鲁宾、赎罪盖，以及存放玛纳——那天使之粮——的金罐。这些大概是用埃及人最精的金制成；次等的金或许制了靠近内幔的纯金灯台及其灯盏，以及放置供饼的金桌，和这二者之间的金香坛。若有第三、第四等的金，则制成圣器。埃及的银亦制成他物；因为以色列子民从寄居埃及中获得极大益处，即得如此大量宝贵材料以供天主服务之用。埃及的衣裳大概制成了圣经中提到的所有刺绣所需；刺绣工以天主的智慧工作，制成如此目的之衣裳，以造幔子及内外院。此刻不宜详论此主题，亦不宜显示以色列子民从埃及人所取之物在多少方面有用。埃及人未曾善用；但希伯来人用以宗教目的，因天主的智慧与他们同在。然而圣经知道，从以色列子民之地下到埃及是恶事；其中蕴含一大真理。因为有些人，在登入天主的律法及以色列的敬拜之后，再与埃及人——即世俗学问——同住，是为恶事。\Person{厄东人阿德尔}，当他在\Place{以色列}地、未尝埃及之饼时，未制造偶像；但当他逃离智慧的\Person{撒罗满}、下到\Place{埃及}，作为逃离天主智慧之人，他与\Person{法郎}联结，娶其妻之妹，生子在法郎众子中长大。因此，虽他回到\Place{以色列}地，却是回来使天主之民分裂，并使他们向金牛犊说：\BibleQuote{以色列，这些就是你的天主，领你出埃及地的神。}我由经验得知，且能告诉您，从埃及取得有用之物而出、再为天主服务预备所需者寥寥无几；而\Person{厄东人阿德尔}却有许多弟兄。我指那些基于某些希腊学问而产出异端思想，如同在\Place{贝特耳}——即天主之家——制造金牛犊的人。据我理解，这是要传达：此等人在圣经中设立自己的形象，而圣经乃天主圣言居所，故寓意上称为\Place{贝特耳}。经文说另一像设在\Place{丹}。\Place{丹}的边界在极处，毗邻异教之地，如\BibleBook{若苏厄}{书}所明载。这些像中有些靠近异教边界，是由\Person{阿德尔}的弟兄们（如我们所指出的）所设计的。\par

% structure: p0008 source=h2 target=section reason=relative-heading-rank; base=h2

\section{三、个人恳请}

那么，先生，我儿，你首先要研读圣经。我说，研读它们。因为我们需要深入地研读圣经著作，以免我们谈论它们比思考更快；当你带着信德和悦乐天主的意向研读这些神圣著作时，要叩击其中封闭之处，守门者必会为你开启，主耶稣论及他说：\BibleRef{若 10:3} \BibleQuote{守门者就给他开门。} 当你专注于这种神圣阅读时，要以正直和坚定不移的信德寻求那存在于圣经多数段落中的隐秘意义。不要满足于叩门和寻求，因为理解神圣事物最需要的是祈祷，救主敦促我们这样做时，不仅说：\BibleRef{玛 7:7} \BibleQuote{叩门，就给你们开门，} 和 \BibleQuote{寻求，就必寻见，} 而且又说：\BibleQuote{求，就必得着。} 我因着对你为父的爱，斗胆说了这些。我做得是否妥当，天主知道，他的基督知道，以及那有分于天主的圣神和基督的圣神的人知道。愿你有分于这些；愿你在这其中的份不断增长，好使你能不仅说：\BibleQuote{我们有分于基督，} \BibleRef{希 3:14} 而且说：\Quote{我们有分于天主。}\par

\SourceRecord{Translated by Allan Menzies.；Ante-Nicene Fathers；Vol. 9.；Edited by Allan Menzies.；Buffalo, NY: Christian Literature Publishing Co.,；1896.；<http://www.newadvent.org/fathers/1014.htm>.}
